% Clase 1 — Agentes Lógicos: Inferencia Proposicional
% Lectura: AIMA Cap. 7
\section{Agentes Lógicos II: Inferencia Proposicional}

\subsection{Reglas de Inferencia}

\begin{tabular}{ll}
\toprule
\textbf{Regla} & \textbf{Forma} \\
\midrule
Modus Ponens & $\alpha \Rightarrow \beta,\ \alpha \vdash \beta$ \\
Modus Tollens & $\alpha \Rightarrow \beta,\ \neg\beta \vdash \neg\alpha$ \\
Resolución & $(\ell_1 \lor \cdots \lor \ell_k),\ (\neg\ell_1 \lor m_1 \lor \cdots) \vdash (m_1 \lor \cdots)$ \\
\bottomrule
\end{tabular}

\subsection{Resolución}

Algoritmo de inferencia completo para lógica proposicional en CNF:
\begin{enumerate}
    \item Convertir KB $\land \neg\alpha$ a CNF.
    \item Aplicar resolución hasta obtener la cláusula vacía (contradicción) o no poder seguir.
    \item Si se obtiene contradicción $\Rightarrow$ KB $\models \alpha$.
\end{enumerate}

\subsection{DPLL y WalkSAT}

\begin{itemize}
    \item \textbf{DPLL}: backtracking completo sobre asignaciones de verdad, con propagación de cláusulas unitarias.
    \item \textbf{WalkSAT}: búsqueda local incompleta, eficaz en la práctica para instancias satisfacibles.
\end{itemize}

\subsection{Ejemplo: Refutación por Resolución}

\begin{example}[Refutación]
    Sea $KB = \{P,\ P \Rightarrow Q,\ Q \Rightarrow R\}$ y se quiere probar $KB \models R$. Se convierte $KB \land \neg R$
    a CNF: $\{P,\ \neg P \lor Q,\ \neg Q \lor R,\ \neg R\}$, y se aplica resolución hasta obtener la cláusula vacía
    $\square$ (una contradicción), lo que confirma el entailment.
\end{example}

\begin{figure}[H]
\centering
\begin{tikzpicture}[
  clause/.style={draw=black, thick, rectangle, rounded corners=1pt, inner sep=4pt, font=\small}
]
\node[clause] (P)   at (0,0)      {$P$};
\node[clause] (PQ)  at (2.6,0)    {$\neg P \lor Q$};
\node[clause] (Q)   at (1.3,-1.4) {$Q$};

\node[clause] (QR)  at (4.2,-1.4) {$\neg Q \lor R$};
\node[clause] (R)   at (2.75,-2.8) {$R$};

\node[clause] (NR)  at (5.5,-2.8) {$\neg R$};
\node[clause] (NIL) at (4.1,-4.2) {$\square$};

\draw[-{Latex},thick] (P)  -- (Q);
\draw[-{Latex},thick] (PQ) -- (Q);
\draw[-{Latex},thick] (Q)  -- (R);
\draw[-{Latex},thick] (QR) -- (R);
\draw[-{Latex},thick] (R)  -- (NIL);
\draw[-{Latex},thick] (NR) -- (NIL);
\end{tikzpicture}
\caption{Árbol de refutación por resolución: cada nodo se obtiene combinando los dos que apuntan hacia él; llegar a la
cláusula vacía $\square$ demuestra que $KB \land \neg R$ es insatisfacible, es decir, $KB \models R$.}
\end{figure}

\begin{exercise}
    Sea $KB = \{A \lor B,\ \neg A \lor C,\ \neg B \lor C,\ \neg C \lor D\}$. Se quiere probar $KB \models D$ por
    refutación.
    \begin{enumerate}
        \item Escriba la cláusula que se agrega a $KB$ al negar la meta.
        \item Aplique resolución sobre $A \lor B$ y $\neg A \lor C$, y por separado sobre $\neg B \lor C$, para derivar
        una cláusula con únicamente $C$.
        \item Combine el resultado anterior con $\neg C \lor D$ y con la meta negada para llegar a la cláusula vacía.
        \item Dibuje el árbol de refutación completo, análogo al de la Figura anterior.
    \end{enumerate}
\end{exercise}

